%=============================================================================!
% 9.6  STABILITY AND THE SIZE OF THE STEP                                     !
%=============================================================================!
\section{Stability and the size of the step}
\label{sec:in-stability}

Chapters~1, 2 and~4 promised to say how long a step molecular dynamics
can take. Two limits answer it: a step beyond which the computed motion
grows without bound, and, well inside it, the step at which the energy
fluctuates by an acceptable amount. This section derives the first,
measures the second, and turns both into a procedure.

\subsection*{The stability limit}

The shadow energy of \cref{eq:in-shadow} contains the factor $1 -
\omega^2\dt^2/4$. For $\omega\dt < 2$ it is positive, $\tilde{H}$ is a sum
of two positive squares, and its curves are ellipses: the computed motion
stays on one and is bounded. For $\omega\dt > 2$ the factor is negative,
and $\tilde{H} = p^2/(2m) - \frac12 kq^2(\omega^2\dt^2/4 - 1)$ is a
difference of two squares, which stays fixed while $p$ and $q$ both grow
without limit. Its curves are hyperbolas, open curves that run off to
infinity, and the computed motion follows one outwards, its size
multiplied, after the first few steps, by a fixed factor at every step
(\cref{ex:in-growth}). Velocity Verlet is therefore stable for the spring
only when
\begin{equation}
   \omega\,\dt < 2 ,
   \label{eq:in-stable}
\end{equation}
that is, with more than $2\pi/2 = \pi$ steps per period.
\Cref{fig:in-stability} measures it: in 200 steps $\lvert q\rvert$ stays
at 1 up to $\omega\dt = 1.99$, reaches $\num{1.2e17}$ at $2.01$ and
$\num{3e38}$ at $2.05$.

\begin{figure}[htb]
\centering
\includegraphics{ch09_integrators/stability}
\caption{The largest $\lvert q\rvert$ reached in 200 steps of velocity
Verlet for a spring started at $q = 1$ at rest, against $\omega\dt$, on a
logarithmic axis; the limit $\omega\dt = 2$ is dotted.}
\label{fig:in-stability}
\end{figure}

Atoms near a minimum of $U$ vibrate in normal modes (\cref{sec:os-modes}).
We write the displacements from the minimum as $\vb{q} = \sum_k
c_k\vb{u}_k$, with $\vb{u}_k$ the patterns of \cref{sec:os-modes}, which
satisfy $\boldsymbol{\Phi}\vb{u}_k = \omega_k^2\vb{M}\vb{u}_k$. When the
forces are linear, pattern $\vb{u}_k$ feels the force
$-\boldsymbol{\Phi}\vb{u}_k = -\omega_k^2\vb{M}\vb{u}_k$, so its
acceleration is $-\omega_k^2\vb{u}_k$: a half kick changes the rate of
change of each $c_k$ by $-\frac12\dt\,\omega_k^2c_k$, a drift changes each
$c_k$ by $\dt$ times its own rate, and velocity Verlet steps every $c_k$ as
it steps a spring of angular frequency $\omega_k$. Each mode must
therefore satisfy \cref{eq:in-stable}: the fastest vibration in the system
sets the limit, however slow everything else is. For the faster O--H
stretch of water, the antisymmetric one, with a period of
\SI{8.88}{\femto\second} (\cref{sec:os-molecules}), $\omega = 2\pi/8.88 =
\SI{0.708}{\radian\per\femto\second}$, and the limit is $\dt < 2/\omega =
\SI{2.83}{\femto\second}$.

\subsection*{Accuracy}

Inside the limit, the energy band of \cref{sec:in-shadow} sets how long a
step is acceptable: for a vibration of angular frequency $\omega$ it
fluctuates by a fraction $\omega^2\dt^2/4$ of the vibration's energy. For
the O--H stretch, a step of \SI{0.5}{\femto\second} gives $\omega\dt =
0.354$ and a band of 3.1\%; \SI{1}{\femto\second} gives 12.5\%. These are
fractions of the energy of the fastest vibration only, but they show why a
step must be short compared with $1/\omega = \SI{1.41}{\femto\second}$, as
\cref{sec:os-molecules} anticipated.

\Cref{fig:in-water} measures this on the spring model of water of
\cref{sec:ro-freedom}, whose fastest vibration, the symmetric stretch,
which the model puts 24\% above the measured value, has $\omega^2 =
\SI{0.7338}{\per\femto\second\squared}$, so $\omega =
\SI{0.857}{\radian\per\femto\second}$ and $2/\omega =
\SI{2.33}{\femto\second}$. Given \SI{0.1}{\electronvolt} of kinetic
energy and followed for \SI{2000}{\femto\second}, its total energy spreads
by \SI{1.3e-3}{\electronvolt}, 1.3\% of the energy put in, with $\dt =
\SI{0.5}{\femto\second}$, and by 9.6\% with \SI{1}{\femto\second}. The
spread grows as $\dt^2$ for the shorter steps (fitted slope 2.02 from 0.05
to \SI{0.5}{\femto\second}) and faster beyond, by a factor of 7.5 from 0.5
to \SI{1}{\femto\second}, where $\omega\dt = 0.86$ is no longer small. The
run blows up at $\dt = \SI{2.3}{\femto\second}$, just inside the limit for
a perfect spring. There $1 - \omega^2\dt^2/4 = 0.030$, so by
\cref{eq:in-shadow} the computed swing of the fastest vibration is
$1/\sqrt{0.030} = 5.8$ times the true one, and over so wide a swing the
springs between atoms are no longer linear in the positions; with
\SI{0.01}{\electronvolt} the same step stays bounded.

\begin{figure}[htb]
\centering
\includegraphics{ch09_integrators/water}
\caption{The spring model of water (\cref{sec:ro-freedom}) followed by
velocity Verlet from its equilibrium geometry with
\SI{0.1}{\electronvolt} of kinetic energy. (a) The total energy relative
to its initial value $E_0 = \SI{0.1}{\electronvolt}$, over the first \SI{100}{\femto\second} for three
steps. (b) The spread of the total energy over \SI{2000}{\femto\second}
against $\dt$, with a line of slope 2 (dashed) and the stability limit
$2/\omega$ of its fastest vibration (dotted); the run at
\SI{2.3}{\femto\second} blew up and is not shown.}
\label{fig:in-water}
\end{figure}

\begin{practice}
The step is chosen by measurement. Without a thermostat, the device of
Chapter~13 that holds the temperature, the total energy should stay
constant, and short runs at two or three steps, such as 0.25, 0.5 and
\SI{1}{\femto\second}, each covering many periods of the vibrations,
record how well it does. Its fluctuation should fall close to $\dt^2$
between the shorter steps, as expected for smooth forces consistent with
the reported energy; it should show no steady
drift; and it should be small compared with the fluctuation of the kinetic
energy, the amount that passes back and forth between kinetic and
potential energy. For the water of \cref{fig:in-water} the ratio of their
standard deviations is 0.9\% at \SI{0.5}{\femto\second} and 6.7\% at
\SI{1}{\femto\second}; how small is acceptable depends on what is to be
measured, and is a judgement. The TrajCast paper gives 0.5 to
\SI{1.0}{\femto\second} as the typical step of molecular dynamics with
forces from electronic structure\cite{Thiemann2026}, and the band of the
O--H stretch above, 3.1\% at \SI{0.5}{\femto\second} and 12.5\% at
\SI{1}{\femto\second}, shows why bonds to hydrogen hold it there. Without
hydrogen the fastest vibrations are slower, and longer steps pass the same
test (for the lithium model alone, \cref{ex:in-lithium}; in graphite the
vibrations of the carbon atoms must pass it too). At half the limit,
$\omega\dt = 1$, the band of the fastest vibration is already 25\%.
Chapter~11 removes the fastest vibrations by holding bonds at fixed
length.
\end{practice}

The stability limit has a cousin in \cref{sec:mo-atoms}: saving the
positions of an oscillation needs enough frames per period to recover its
velocity, and computing the motion needs more than $\pi$ steps per period
to remain bounded at all.

\begin{keyidea}
Velocity Verlet is stable only for $\omega\dt < 2$ for every vibration, so
the fastest one sets the limit, \SI{2.83}{\femto\second} for the O--H
stretch. Accuracy demands a much shorter step: the energy of a vibration
fluctuates by $\omega^2\dt^2/4$, and the step is chosen by measuring the
energy fluctuation of short runs.
\end{keyidea}

\notebook{09}{push the step to the stability limit}

\begin{selfcheck}
The model lithium atom rattling in its hollow (\cref{sec:os-small}) has a
period of \SI{170.3}{\femto\second}. What is the stability limit for the
step?
\end{selfcheck}
